\section{参考项目与总体设计}
\subsection{参考项目}
酱菜通面向江西财经大学学生，主要任务是在一个客户端内组织课表、学习任务、个人待办与生活查询。页面和数据整理采用本地 WebView 与 JavaScript，登录会话、受控网络请求及系统能力由 Java 原生层处理。已有组件与目标校园的实际条件决定了各模块的实现范围。江财校历和课表以校方数据为准，待办可在离线环境完整使用，各外部服务仍分别授权。

航院通为移动端架构提供参考。酱菜通沿用其中的 WebView 与 Java 桥接方式，复用了学习通的部分查询、解析、扫描与排序代码\cite{hangyuan}。慧生活798以 life-798 为主要技术参考，用于理解设备账号与积分账号的分工、验证码流程与会话管理，并核对余额、设备、用水控制、账单和积分处理\cite{life798}。排查真实登录问题时，另以已验证可登录的客户端核对请求及会话行为，具体过程见第四章。

\subsection{酱菜通的设计与实现}
江财教务的学期选项、周次表达和成绩报表需要依据实际校方页面处理。酱菜通在同一课程模型中连接校历、课程时段与缓存，并集中提供教务会话的授权及重连入口。学习任务使用教学班组织，活跃度规则另外判断课程是否在默认列表中显示，避免把历史课程和重要任务一并过滤。个人待办用于记录平台之外的日常事务，数据写入本机后可以独立于各服务的登录状态使用。生活页与账号中心分别承接业务操作和授权反馈，住理生活按当前接入状态提供官方服务入口。表\ref{tab:compare}列出这些设计及相关参考。

\begin{table}[htbp]
\centering\caption{酱菜通主要业务及技术参考}\label{tab:compare}
\small
\begin{tabularx}{\linewidth}{@{}>{\raggedright\arraybackslash}p{23mm}YY@{}}
\toprule
业务领域 & 酱菜通设计与实现 & 借鉴的技术思路 \\
\midrule
移动端架构 & 区分应用主界面、校方认证和本机存储的职责 & WebView 与 Java 桥接的端侧组织 \\
江财教务 & 校方课表读取、校历映射、有效成绩与考试报表查询 & 校园服务的模块分组 \\
学习任务 & 同轮处理教学班分项任务与隐藏课程复查，按账号保存可见偏好 & 分项读取、解析、扫描与任务排序 \\
校园生活 & 分别管理设备和积分账号，提供用水查询，并依接入状态展示洗浴入口 & life-798 的双账号分工、验证码流程和业务请求 \\
待办与设置 & 本地保存待办，集中处理账号连接与重连，统一页面层级 & 服务账号和业务页面的入口组织 \\
\bottomrule
\end{tabularx}
\end{table}

\subsection{系统总体架构}
图\ref{fig:arch}展示从页面操作到服务返回的两条数据路径。主界面由 Android 承载本地 HTML、CSS 与 JavaScript。受信的应用页面把请求参数交给 Java 桥接，原生层管理请求范围和对应会话，再将结果送回页面的回调。学习通和798沿用这一路径，JavaScript 负责把响应整理成页面可使用的条目，再按业务状态刷新列表。需在网页内完成的校方认证及页面读取走独立 WebView，院校数据返回主界面前再核验账号、学期和结构。本项目把主界面和外部网页分开，避免面向应用资产的原生能力被注入到认证网页中。

本机保存内容按用途安排。课表及学习结果首先服务于回读，课表按账号与学期隔离，学习缓存按账号隔离。待办和主题偏好用本机 WebView 存储。需要保存的原生凭据由应用私有存储承载，使用 Android Keystore 中的密钥进行加密，代码采用 AES/GCM。凭据留在原生层管理，页面只需获得完成查询所需的返回数据和连接状态。将业务缓存、用户自己录入的待办和会话分别保存，也使普通网页缓存清理可以保留账号及待办。

\begin{figure}[htbp]
\centering\begin{tikzpicture}
\node[block,text width=91mm] (web) at (0,0) {应用主界面 WebView\\操作、业务模型与结果展示};
\node[block,text width=43mm] (bridge) at (-4.05,-1.8) {Java 原生桥接\\受控请求与结果回调};
\node[block,text width=43mm] (portal) at (4.05,-1.8) {独立教务 WebView\\主动认证与后台页面读取};
\node[block,text width=42mm] (store) at (0,-3.6) {本机存储\\缓存、待办和原生凭据};
\node[block,text width=35mm] (vendor) at (-5.35,-3.6) {学习通／798};
\node[block,text width=35mm] (school) at (5.35,-3.6) {江财教务门户};
\draw[flow,<->] (web.south -| bridge.north) -- (bridge.north);
\draw[flow,<->] (web.south -| portal.north) -- (portal.north);
\draw[flow,<->] (web.south) -- node[note,right,pos=.63] {业务数据} (store.north);
\draw[flow,<->] (bridge.south) -- ++(0,-3mm) -| (vendor.north);
\draw[flow,<->] (portal.south) -- ++(0,-3mm) -| (school.north);
\draw[flow,<->] (bridge.east) -| node[note,above,pos=.27] {凭据} ([xshift=-10mm]store.north);
\end{tikzpicture}

\caption{本地页面、原生桥接和外部服务的数据流}\label{fig:arch}
\end{figure}

\subsection{功能模块组织}
功能结构见图\ref{fig:modules}。底栏按任务场景安排，账号管理集中在“我的”。教务、学习通、798设备账号、798积分账号和住理生活在这里显示各自的连接状态。业务页负责查询和查看结果，设置页承接用户明确发起的绑定或重连。这个划分在课表失效时尤其重要，课表仍可保留原有内容，用户自行决定何时重新认证。

\begin{figure}[htbp]
\centering\begin{tikzpicture}
\node[block,text width=46mm] (app) at (0,0) {酱菜通};
\foreach \x/\name/\content in {
  -5.86/a/{教务\\课表\\成绩与考试},
  -2.93/b/{学习\\任务筛选\\课程活跃度},
  0/c/{待办\\编辑与完成\\离线保存},
  2.93/d/{生活\\798 饮水\\住理生活},
  5.86/e/{我的\\账号管理\\外观与缓存}
}{
  \node[block,text width=24mm,minimum height=22mm,font=\footnotesize] (\name) at (\x,-2.1) {\content};
  \draw[flow] (app.south) -- (0,-1.0) -| (\name.north);
}
\end{tikzpicture}

\caption{五个业务模块及其功能分组}\label{fig:modules}
\end{figure}
